% ============================================================================
% Chapter 7: Discussion, Conclusion & Future Work
% ============================================================================
\chapter{Discussion, Conclusion, and Future Work}
\label{ch:conclusion}

\section{Discussion and Limitations}

\subsection{Transition to Semi-Supervised Operation}

CBSR's Behavioral Pattern Reasoner and meta-learner are calibrated using sparse validation labels,
transitioning the framework from fully unsupervised to \textbf{semi-supervised}. This is a
deliberate design trade-off: while a small validation overhead is introduced, the BPR learns an
explicit decision boundary between benign noise and attack patterns, insulating the system from the
calibration poisoning vulnerability that afflicted the pure causal baseline (v1), and providing far
higher classification performance. The state representation itself (causal discovery, SCM fitting,
and state projection) remains fully unsupervised and requires no labels.

\subsection{Calibration Sensitivity Under Contamination}

Under the pure causal baseline (v1), attack events leaking into the conformal calibration queue
inflated the empirical-CDF baseline, masking subsequent test attacks and causing AUC to drop from
0.972 to 0.681 under just 1\% training contamination. CBSR significantly mitigates this through
RobustParCorr (MCD-based, tolerating 15\% contamination in causal discovery) and the BPR's
explicit label-supervised training. However, very high contamination rates ($>$20\%) remain an
open challenge.

\subsection{Scalability Limitations on High-Dimensional Logs}

PCMCI causal discovery on BETH ($d = 882$ features, $T = 360$ training windows) requires
7,842 seconds ($\approx$2.18 hours) on a standard CPU. This is a one-time offline or
asynchronous background process. For environments with thousands of unique edge types ($d > 1000$),
a mandatory feature selection pre-pass (variance thresholding or mutual-information ranking) is
required before PCMCI. GPU-accelerated robust covariance estimation and distributed PCMCI
are targeted as engineering extensions.

\subsection{Causal Sufficiency Assumption}

CBSR assumes causal sufficiency---no unobserved confounders jointly drive the observed provenance
features. In enterprise networks, unobserved external variables (e.g., domain controller
group-policy pushes, NTP synchronization) could induce spurious causal edges, resulting in
localized false alarms. Two mechanisms mitigate this: (1) the conformal prediction layer absorbs
Type I error inflation from spurious edges, and (2) the change-point detector triggers periodic
online refits that discard stale or spurious edges. The practical consequence is elevated (but
bounded) false alarms, not missed attacks---a favorable failure mode for a defense system.

\subsection{StreamSpot and Heterogeneous Provenance}

StreamSpot (AUC 0.6909) highlights a clear boundary condition: a single linear SCM learned from
a mixture of heterogeneous workload types is a poor fit for any individual type. The temporal
reasoner ($R_T$) rescues much of the performance by capturing sequential attack signatures, but
the mechanism violation coordinates remain desensitized. Workload-conditioned SCMs or
mixture-of-experts architectures are needed for production deployment on heterogeneous graph
databases.

\subsection{Mimicry Attacks}

Successful SCM-aware mimicry requires an attacker to simultaneously satisfy all $d \times
\tau_{\max}$ causal constraints while executing a multi-stage attack---a highly constrained
optimization problem. Moreover, real APT campaigns inject novel process--file edges with no
causal parents in the baseline graph, triggering the structural novelty coordinate ($s_1$)
regardless of residual behavior. Empirical evaluation of active mimicry adversaries is left to
future work.


\section{Answering the Research Questions}

\textbf{RQ1: Online causal discovery from unlabeled data.}
CBSR demonstrates that RobustParCorr-based PCMCI can discover stable causal mechanisms from
streaming, unlabeled, and potentially contaminated provenance data. The MCD estimator tolerates
up to 15\% contamination in causal structure learning, confirmed by the contamination sweep in
Section~\ref{sec:contamination}.

\textbf{RQ2: Rich behavioral state representation for multi-perspective reasoning.}
The 12-dimensional Behavioral State Space successfully captures diverse APT manifestations across
five coordinate groups. Feature importance analysis confirms that Mechanism Violation coordinates
($s_9$: Log Min $p$-value, $s_{10}$: Residual Energy) are the primary detection drivers, while
the Causal Intervention Score ($s_{12}$) provides the strongest topological signal on complex
enterprise APTs (OpTC importance = 0.332).

\textbf{RQ3: Online threshold calibration under drift.}
The conformal prediction feedback loop maintains empirical false alarm rates at 1.78\%--6.13\%
across all datasets at a 5\% FPR budget. The change-point-triggered SCM refit and queue flush
hold post-drift false alarm rates to 9.2\%---3.2$\times$--5.6$\times$ lower than static
alternatives.


\section{Summary of Contributions}

This thesis presented \textbf{Causal Behavioral State Reasoning (CBSR)}, a novel
provenance-based IDS that eliminates the clean-data and label assumptions that constrain all
prior work in this domain. The key contributions are:

\begin{enumerate}
\item \textbf{Formal Behavioral State Space}: A 12-dimensional state representation grounded in
    five invariant coordinate groups derived from online SCM outputs, providing a principled
    basis for multi-perspective threat reasoning on streaming provenance telemetry.

\item \textbf{Causal Intervention Score (CIS)}: A graph-boundary metric grounded in spectral
    graph theory that quantifies multi-hop violation propagation, separating localized
    administrative noise from cascading APT lateral movement.

\item \textbf{Multi-Perspective Parallel Reasoning}: A three-operator architecture ($R_L$,
    $R_T$, $R_B$) that resolves the model-mismatch failures of pure causal detectors by
    combining instantaneous, temporal, and behavioral perspectives on the shared state space.

\item \textbf{Calibrated Risk Inference with Conformal Guarantees}: A stacked meta-fusion layer
    with online conformal prediction providing distribution-free FPR bounds under concept drift,
    with change-point-triggered queue flushing to restore localized exchangeability.

\item \textbf{Multi-Benchmark Validation}: AUC of 0.9628 on DARPA OpTC (real enterprise Windows
    logs), AUC 0.9901 on BETH (Linux sysdig telemetry), AUC 0.9837 / 0.9784 on TC3 / NODLINK
    (simulated APT campaigns), and AUC 0.6909 on StreamSpot---a $+$19.2\% improvement over the
    pure causal baseline on this challenging heterogeneous benchmark.
\end{enumerate}


\section{Future Work}

We identify four key areas for future exploration:

\begin{enumerate}
\item \textbf{Workload-Conditioned SCM Ensembles}: The StreamSpot weakness stems from fitting a
    single linear SCM over heterogeneous graph workloads. Future work will explore workload-type
    conditioning via graph clustering pre-passes or mixture-of-experts SCMs with soft workload
    assignment, targeting substantially higher detection on heterogeneous-database environments.

\item \textbf{Scalable Causal Discovery for High-Dimensional Logs}: For $d > 500$ edge types,
    PCMCI discovery time becomes a practical bottleneck. Future work will integrate
    GPU-accelerated FastMCD robust covariance estimation and distributed PCMCI to reduce offline
    startup time to under 6 hours for $d = 5000$ features.

\item \textbf{Robust Conformal Calibration under Heavy Contamination}: Developing trimmed eCDF
    estimation or self-supervised anomaly filtering to screen the rolling conformal calibration
    queue, improving resilience when training streams carry $>$20\% contamination.

\item \textbf{Empirical SCM-Aware Mimicry Attack Evaluation}: Designing and empirically
    evaluating adversarial scenarios where an attacker actively learns and mimics the discovered
    SCM structure, and evaluating whether the structural novelty coordinate and CIS can maintain
    robustness against this sophisticated threat.
\end{enumerate}


\section{Closing Remarks}

CBSR presents a paradigm shift in provenance-based IDS design. By projecting raw system telemetry
into a formal causal behavioral state space and analyzing it from multiple reasoning perspectives,
it opens a practical pathway for enterprise deployments where clean training data is unavailable.
The formal state representation, CIS metric, and multi-perspective reasoning architecture together
provide a strong foundation for future research at the intersection of causal inference and
cybersecurity.
